% Candidate only.  The cited source statements and the live-book interfaces
% are listed explicitly in SOURCE_LEDGER.md.
\begin{proof}
By the invariant-function correspondence (5.5) of \cite{BGJKK25}, there
are unique finite continuous functions $\varphi,\psi$ on $O$ such that
\[
 u=\varphi\circ\operatorname{trop},\qquad
 v=\psi\circ\operatorname{trop}.
\]
The test-form computation in the proof of Proposition~5.3.10(iv) of
\cite{BGJKK25} uses only this factorization and continuity: pullbacks of
compactly supported top-degree tropical forms are supported on the canonical
skeleton.  It therefore gives, for the two present continuous functions,
\[
 [\varphi]=\operatorname{trop}_*[u],\qquad
 [\psi]=\operatorname{trop}_*[v].
\]
The same proof displays commutation of this pushforward with
$\dprimeddprime$.  Thus
\[
 \dprimeddprime[\varphi]
 =\operatorname{trop}_*\dprimeddprime[u]
 =\operatorname{trop}_*\dprimeddprime[v]
 =\dprimeddprime[\psi].
\]

Put $h=\varphi-\psi$.  On every Euclidean ball $B\Subset O$, the equality
$\dprimeddprime[h]=0$ says that every distributional second partial
derivative of $h|_B$ is zero.  Hence $h|_B$ is affine.  (Indeed each first
distributional derivative is constant on the connected ball; the continuous
function $h$ therefore equals the resulting affine function.)  Thus
$h\circ\operatorname{trop}$ is locally a smooth CLD function with zero
$\dprimeddprime$: after choosing a basis of $N$, it is locally an affine
function of the moment coordinates $-\log|\chi^m|$; in the book's notation
this is \cref{lem:ch19-toric-chart}.

Now $u-v=h\circ\operatorname{trop}$ is smooth and has zero CLD Hessian.
The closed-current Leibniz identity, equivalently
\cref{cor:ch16-smooth-h}, gives
\[
 \dprimeddprime\bigl((u-v)T\bigr)=0.
\]
Linearity of multiplication by $T$ proves the asserted equality.
\end{proof}
